\documentclass[11pt]{article}
% Round-2 letter. Kept for the record: its PDF is frozen; recompiling it would pick up later numbers.
\usepackage[T1]{fontenc}
\usepackage[a4paper,margin=2.2cm]{geometry}
\usepackage{amsmath,amssymb}
\usepackage[per-mode=symbol]{siunitx}
\usepackage{enumitem}
\usepackage{xcolor}
\usepackage[hidelinks]{hyperref}
% Auto-generated by paper/make_paper_assets.py - do not edit by hand.

\newcommand{\YieldL}{1608}
\newcommand{\PRL}{0.819}
\newcommand{\CUFL}{18.35}
\newcommand{\POAL}{1963}
\newcommand{\TdayL}{39.5}
\newcommand{\TmaxL}{57.3}
\newcommand{\TriseL}{10.1}
\newcommand{\RHL}{74.0}
\newcommand{\DewL}{1{,}163}
\newcommand{\SRL}{0.965}
\newcommand{\DegBaseL}{0.30}
\newcommand{\DegTHL}{0.34}
\newcommand{\DegPIDL}{0.15}
\newcommand{\DegCorrL}{0.09}
\newcommand{\DegTotL}{0.88}
\newcommand{\TOWL}{3{,}055}
\newcommand{\RHencL}{63.5}
\newcommand{\LifeL}{36{,}233}
\newcommand{\CapL}{80.9}
\newcommand{\LCOEL}{2.83}
\newcommand{\LossIncidenceAngleL}{3.44}
\newcommand{\LossSpectralL}{-0.23}
\newcommand{\LossSoilingL}{3.67}
\newcommand{\LossDewCondensationL}{0.02}
\newcommand{\LossModuleTemperatureL}{6.12}
\newcommand{\LossDcWiringMismatchMotionLidL}{3.46}
\newcommand{\LossInverterAndClippingL}{2.03}
\newcommand{\LossAvailabilityL}{1.00}
\newcommand{\YieldF}{1634}
\newcommand{\PRF}{0.837}
\newcommand{\CUFF}{18.66}
\newcommand{\POAF}{1953}
\newcommand{\TdayF}{36.5}
\newcommand{\TmaxF}{50.0}
\newcommand{\TriseF}{7.3}
\newcommand{\RHF}{77.5}
\newcommand{\DewF}{1{,}092}
\newcommand{\SRF}{0.988}
\newcommand{\DegBaseF}{0.30}
\newcommand{\DegTHF}{0.34}
\newcommand{\DegPIDF}{0.15}
\newcommand{\DegCorrF}{0.05}
\newcommand{\DegTotF}{0.84}
\newcommand{\TOWF}{3{,}719}
\newcommand{\RHencF}{68.4}
\newcommand{\LifeF}{37{,}004}
\newcommand{\CapF}{81.7}
\newcommand{\LCOEF}{3.54}
\newcommand{\LossIncidenceAngleF}{3.47}
\newcommand{\LossSpectralF}{-0.24}
\newcommand{\LossSoilingF}{1.27}
\newcommand{\LossDewCondensationF}{0.04}
\newcommand{\LossModuleTemperatureF}{4.78}
\newcommand{\LossDcWiringMismatchMotionLidF}{4.62}
\newcommand{\LossInverterAndClippingF}{2.03}
\newcommand{\LossAvailabilityF}{1.50}
\newcommand{\YieldS}{1590}
\newcommand{\PRS}{0.814}
\newcommand{\CUFS}{18.15}
\newcommand{\POAS}{1953}
\newcommand{\TdayS}{36.2}
\newcommand{\TmaxS}{49.4}
\newcommand{\TriseS}{7.0}
\newcommand{\RHS}{79.7}
\newcommand{\DewS}{1{,}389}
\newcommand{\SRS}{0.969}
\newcommand{\DegBaseS}{0.30}
\newcommand{\DegTHS}{0.36}
\newcommand{\DegPIDS}{0.21}
\newcommand{\DegCorrS}{0.30}
\newcommand{\DegTotS}{1.16}
\newcommand{\TOWS}{4{,}288}
\newcommand{\RHencS}{70.4}
\newcommand{\LifeS}{34{,}677}
\newcommand{\CapS}{75.5}
\newcommand{\LCOES}{4.43}
\newcommand{\LossIncidenceAngleS}{3.47}
\newcommand{\LossSpectralS}{-0.25}
\newcommand{\LossSoilingS}{3.29}
\newcommand{\LossDewCondensationS}{0.05}
\newcommand{\LossModuleTemperatureS}{4.64}
\newcommand{\LossDcWiringMismatchMotionLidS}{4.91}
\newcommand{\LossInverterAndClippingS}{2.03}
\newcommand{\LossAvailabilityS}{2.00}
\newcommand{\YieldLM}{1608}
\newcommand{\PRLM}{0.819}
\newcommand{\CUFLM}{18.35}
\newcommand{\POALM}{1963}
\newcommand{\TdayLM}{39.5}
\newcommand{\TmaxLM}{57.3}
\newcommand{\TriseLM}{10.1}
\newcommand{\RHLM}{74.0}
\newcommand{\DewLM}{1{,}163}
\newcommand{\SRLM}{0.965}
\newcommand{\DegBaseLM}{0.30}
\newcommand{\DegTHLM}{0.34}
\newcommand{\DegPIDLM}{0.02}
\newcommand{\DegCorrLM}{0.04}
\newcommand{\DegTotLM}{0.70}
\newcommand{\TOWLM}{3{,}055}
\newcommand{\RHencLM}{63.5}
\newcommand{\LifeLM}{37{,}002}
\newcommand{\CapLM}{84.5}
\newcommand{\LCOELM}{2.90}
\newcommand{\LossIncidenceAngleLM}{3.44}
\newcommand{\LossSpectralLM}{-0.23}
\newcommand{\LossSoilingLM}{3.67}
\newcommand{\LossDewCondensationLM}{0.02}
\newcommand{\LossModuleTemperatureLM}{6.12}
\newcommand{\LossDcWiringMismatchMotionLidLM}{3.46}
\newcommand{\LossInverterAndClippingLM}{2.03}
\newcommand{\LossAvailabilityLM}{1.00}
\newcommand{\YieldFM}{1634}
\newcommand{\PRFM}{0.837}
\newcommand{\CUFFM}{18.66}
\newcommand{\POAFM}{1953}
\newcommand{\TdayFM}{36.5}
\newcommand{\TmaxFM}{50.0}
\newcommand{\TriseFM}{7.3}
\newcommand{\RHFM}{77.5}
\newcommand{\DewFM}{1{,}092}
\newcommand{\SRFM}{0.988}
\newcommand{\DegBaseFM}{0.30}
\newcommand{\DegTHFM}{0.34}
\newcommand{\DegPIDFM}{0.01}
\newcommand{\DegCorrFM}{0.03}
\newcommand{\DegTotFM}{0.68}
\newcommand{\TOWFM}{3{,}719}
\newcommand{\RHencFM}{68.4}
\newcommand{\LifeFM}{37{,}698}
\newcommand{\CapFM}{84.9}
\newcommand{\LCOEFM}{3.60}
\newcommand{\LossIncidenceAngleFM}{3.47}
\newcommand{\LossSpectralFM}{-0.24}
\newcommand{\LossSoilingFM}{1.27}
\newcommand{\LossDewCondensationFM}{0.04}
\newcommand{\LossModuleTemperatureFM}{4.78}
\newcommand{\LossDcWiringMismatchMotionLidFM}{4.62}
\newcommand{\LossInverterAndClippingFM}{2.03}
\newcommand{\LossAvailabilityFM}{1.50}
\newcommand{\YieldSM}{1590}
\newcommand{\PRSM}{0.814}
\newcommand{\CUFSM}{18.15}
\newcommand{\POASM}{1953}
\newcommand{\TdaySM}{36.2}
\newcommand{\TmaxSM}{49.4}
\newcommand{\TriseSM}{7.0}
\newcommand{\RHSM}{79.7}
\newcommand{\DewSM}{1{,}389}
\newcommand{\SRSM}{0.969}
\newcommand{\DegBaseSM}{0.30}
\newcommand{\DegTHSM}{0.36}
\newcommand{\DegPIDSM}{0.02}
\newcommand{\DegCorrSM}{0.15}
\newcommand{\DegTotSM}{0.83}
\newcommand{\TOWSM}{4{,}288}
\newcommand{\RHencSM}{70.4}
\newcommand{\LifeSM}{36{,}045}
\newcommand{\CapSM}{81.9}
\newcommand{\LCOESM}{4.43}
\newcommand{\LossIncidenceAngleSM}{3.47}
\newcommand{\LossSpectralSM}{-0.25}
\newcommand{\LossSoilingSM}{3.29}
\newcommand{\LossDewCondensationSM}{0.05}
\newcommand{\LossModuleTemperatureSM}{4.64}
\newcommand{\LossDcWiringMismatchMotionLidSM}{4.91}
\newcommand{\LossInverterAndClippingSM}{2.03}
\newcommand{\LossAvailabilitySM}{2.00}
\newcommand{\GainYieldF}{1.6}
\newcommand{\GainYieldS}{-1.1}
\newcommand{\GainLifeF}{2.1}
\newcommand{\GainLifeS}{-4.3}
\newcommand{\GainLifeSM}{3.9}
\newcommand{\dTdayF}{3.0}
\newcommand{\dTmaxF}{7.3}
\newcommand{\WaterSaved}{11{,}700}
\newcommand{\GHI}{1{,}889}
\newcommand{\Tair}{28.1}
\newcommand{\RHair}{74.0}
\newcommand{\Rain}{1{,}043}
\newcommand{\SensLifeMin}{32{,}498}
\newcommand{\SensLifeMax}{37{,}005}
\newcommand{\SensDegMin}{0.78}
\newcommand{\SensDegMax}{1.49}
\newcommand{\SensSaltSpan}{2{,}854}
\newcommand{\SensMoistSpan}{1{,}674}
\newcommand{\SensDegPerK}{0.167}
\newcommand{\SensTOWZero}{2{,}253}
\newcommand{\SensTOWMax}{5{,}547}
\newcommand{\SensYieldLoPct}{4.6}
\newcommand{\BreakEvenSalinity}{$<$0.25}
\newcommand{\NSeeds}{20}
\newcommand{\EnsYieldF}{1.71}
\newcommand{\EnsYieldSdF}{0.08}
\newcommand{\EnsLifeF}{2.17}
\newcommand{\EnsLifeSdF}{0.10}
\newcommand{\EnsDegF}{-0.039}
\newcommand{\EnsDegSdF}{0.003}
\newcommand{\EnsYieldS}{-1.08}
\newcommand{\EnsYieldSdS}{0.05}
\newcommand{\EnsLifeS}{-4.32}
\newcommand{\EnsLifeSdS}{0.06}
\newcommand{\EnsDegS}{+0.288}
\newcommand{\EnsDegSdS}{0.002}
\newcommand{\EnsYieldLandMin}{1{,}570}
\newcommand{\EnsYieldLandMax}{1{,}649}
\newcommand{\EnsFreshWins}{20}

% Auto-generated by paper/make_paper_assets.py (review additions) - do not edit by hand.

\newcommand{\NMC}{1{,}000}
\newcommand{\NParams}{28}
\newcommand{\MCYieldMedF}{1.2}
\newcommand{\MCYieldLoF}{$-$1.7}
\newcommand{\MCYieldHiF}{5.3}
\newcommand{\MCYieldPosF}{78}
\newcommand{\MCLifeMedF}{1.5}
\newcommand{\MCLifeLoF}{$-$3.6}
\newcommand{\MCLifeHiF}{6.3}
\newcommand{\MCLifePosF}{75}
\newcommand{\MCYieldMedS}{$-$1.5}
\newcommand{\MCYieldLoS}{$-$5.5}
\newcommand{\MCYieldHiS}{3.1}
\newcommand{\MCYieldPosS}{27}
\newcommand{\MCLifeMedS}{$-$5.1}
\newcommand{\MCLifeLoS}{$-$11.5}
\newcommand{\MCLifeHiS}{0.7}
\newcommand{\MCLifePosS}{4}
\newcommand{\MCYieldMedSM}{$-$1.5}
\newcommand{\MCYieldLoSM}{$-$5.5}
\newcommand{\MCYieldHiSM}{3.1}
\newcommand{\MCYieldPosSM}{27}
\newcommand{\MCLifeMedSM}{$-$1.6}
\newcommand{\MCLifeLoSM}{$-$7.0}
\newcommand{\MCLifeHiSM}{3.8}
\newcommand{\MCLifePosSM}{27}
\newcommand{\RCNameA}{Land dust deposition}
\newcommand{\RCValA}{0.45}
\newcommand{\RCNameB}{FPV $U_1$}
\newcommand{\RCValB}{0.43}
\newcommand{\RCNameC}{Saline salt deposition}
\newcommand{\RCValC}{$-$0.41}
\newcommand{\RCNameD}{FPV $U_0$}
\newcommand{\RCValD}{0.30}
\newcommand{\RCNameE}{Saline chloride deposition}
\newcommand{\RCValE}{$-$0.28}
\newcommand{\RCNameF}{Land $U_1$}
\newcommand{\RCValF}{$-$0.26}
\newcommand{\RCFNameA}{Land dust deposition}
\newcommand{\RCFValA}{0.59}
\newcommand{\RCFNameB}{FPV $U_1$}
\newcommand{\RCFValB}{0.49}
\newcommand{\RCFNameC}{FPV $U_0$}
\newcommand{\RCFValC}{0.35}
\newcommand{\PLcoeF}{99.5}
\newcommand{\PLcoeS}{100.0}
\newcommand{\OATBase}{$-$4.30}
\newcommand{\ClStarText}{none in 0--400}
\newcommand{\SaltStarYield}{35}
\newcommand{\DTdStar}{2.7}
\newcommand{\LifeSalineClZeroPct}{$-$0.9}
\newcommand{\LifeSalineSaltZeroPct}{$-$0.9}
\newcommand{\YieldSalineSaltZeroPct}{1.6}
\newcommand{\DegSalineClMax}{1.31}
\newcommand{\DegSalineClZero}{0.87}
\newcommand{\ClMax}{400}
\newcommand{\TOWFreshAtStar}{5{,}745}
\newcommand{\DewFreshMin}{1.64}
\newcommand{\DewFreshMax}{1.66}
\newcommand{\DewSalineMin}{$-$1.10}
\newcommand{\DewSalineMax}{$-$1.06}
\newcommand{\MitGGvsLand}{1.5}
\newcommand{\MitAllvsLand}{3.3}
\newcommand{\MitPIDRec}{2.2}
\newcommand{\MitSaltRec}{1.7}
\newcommand{\MitBothRec}{3.9}
\newcommand{\MitBothvsLand}{$-$0.5}
\newcommand{\CapexStarF}{34.3}
\newcommand{\CapexStarS}{29.1}
\newcommand{\LandPremF}{10.5}
\newcommand{\LandPremS}{23.6}
\newcommand{\WaterStar}{93}
\newcommand{\LCOEWaterZero}{3.54}
\newcommand{\LCOEWaterTen}{3.47}
\newcommand{\LCOEWaterTwenty}{3.39}
\newcommand{\LCOEWaterForty}{3.24}
\newcommand{\LCOEWaterEighty}{2.93}
\newcommand{\FinGapFMin}{0.58}
\newcommand{\FinGapFMax}{0.86}
\newcommand{\FinGapSMin}{1.36}
\newcommand{\FinGapSMax}{1.87}

\setlength{\parindent}{0pt}
\setlength{\parskip}{5pt}
\newcounter{cmt}
\newenvironment{comment}[1]{\refstepcounter{cmt}\par\medskip\noindent\textbf{Comment \thecmt{} (#1).}\itshape\ }{\par}
\newenvironment{response}{\par\noindent\textbf{Response.}\ }{\par}
\newcommand{\change}[1]{\par\noindent\textcolor{blue!50!black}{\textbf{Changes:} #1}\par}

\begin{document}
\begin{center}
{\Large\bfseries Response to Reviewers --- Second Round}\\[4pt]
Manuscript: ``Lifetime Performance and Degradation of Floating Photovoltaic Systems Under Humid and Saline Conditions''
\end{center}

I thank the reviewer for the second careful assessment and for the constructive view of the first revision. This round
addresses the remaining concerns as follows:
\begin{enumerate}[nosep]
\item Baseline and Monte Carlo results are now distinguished explicitly in the abstract, results and conclusion.
\item The conditional nature of the freshwater benefit is now a headline conclusion.
\item A full Monte Carlo parameter table with sources has been added (Appendix, Table~XIII).
\item Financing uncertainty is included (discount rate and O\&M escalation, Table~XI and the Monte Carlo).
\item The constant-rate lifetime projection is stated explicitly.
\item The supported physics and the assumed functional forms of the salt--PID coupling are separated.
\item A concrete field-validation protocol has been added (Table~XII).
\end{enumerate}
The two financing parameters raise the Monte Carlo to \NParams{} parameters. Because the random sequence therefore
changes, all Monte Carlo statistics are regenerated. They differ from the first revision only within sampling noise, for
example freshwater median $+$\MCLifeMedF\% instead of $+1.3$\%, and saline median \MCLifeMedS\% instead of $-4.9$\%. The
baseline results are unchanged. As before, every number is generated by one script.

\begin{comment}{\S4, mitigation wording}
Call the mitigation results ``scenario-based mitigation benefits'', not validated benefits.
\end{comment}
\begin{response}
Done. Table~VIII is retitled ``Scenario-Based Mitigation Benefits for the Saline FPV Plant'', and Section~IV-E states
that these are scenario-based mitigation benefits, not validated ones. The ``what-if'' language is retained.
\end{response}

\begin{comment}{\S5, thermal headline}
Make the dependence of the freshwater benefit on the thermal characteristics of the floating structure a headline
conclusion.
\end{comment}
\begin{response}
Done. The abstract now states that ``the freshwater benefit is conditional: it is positive in \MCLifePosF\% of samples,
depends strongly on the thermal characteristics of the floating structure, and vanishes if the dew point over the water
rises by more than about \DTdStar~K.'' The conclusion repeats this, and the new Section~V-B states as a practical rule
that the benefit ``should not be assumed from the presence of water alone''.
\end{response}

\begin{comment}{\S6, weather wording}
Use ``climatologically consistent synthetic weather generator'' rather than ``validated''.
\end{comment}
\begin{response}
Done. Section~III-A now says the generator ``is climatologically consistent with its inputs; it is \emph{not} an
external validation''. The word ``validated'' is not used for the weather model anywhere in the manuscript.
\end{response}

\begin{comment}{\S7 and \S22.6, field validation}
Add a validation case or, if unavailable, present the work as a scenario and uncertainty framework requiring field
calibration, with a concrete validation protocol.
\end{comment}
\begin{response}
Measured FPV data are still not available to this study, so the second option is taken. New Section~V-E and Table~XII
propose a minimum measurement set for an operating FPV plant with a co-located land reference:
\begin{itemize}[nosep]
\item GHI and POA irradiance;
\item air temperature, RH and wind over water and over land;
\item water temperature;
\item module temperature at several positions;
\item a dew/wetness sensor on the glass;
\item cleaned and uncleaned soiling reference cells;
\item chloride deposition by wet candle (ISO~9225);
\item rainfall and string DC power;
\item insulation resistance and leakage current;
\item quarterly IV curves of reference modules.
\end{itemize}
Each measurement is mapped to the model parameter it would calibrate. The abstract and conclusion keep the statement
that the results are ``model-based scenario predictions that still require field validation''.
\end{response}
\change{Section~V-E and Table~XII (new).}

\begin{comment}{\S9 and \S22.5, constant degradation rate}
State that the lifetime projection assumes a constant effective annual degradation rate.
\end{comment}
\begin{response}
Added to Section~III-H, essentially in the reviewer's words: ``The lifetime projection assumes a constant effective
annual degradation rate derived from the stress distribution of the simulated year; the 25-year results therefore
represent an equivalent-rate scenario, not a mechanistic year-by-year degradation trajectory.''
\end{response}

\begin{comment}{\S10, salt--PID coupling}
Distinguish the experimentally supported physics from the assumed leakage function and threshold shift.
\end{comment}
\begin{response}
Done. Section~III-F now separates \emph{experimentally supported physics} (PID leakage rises steeply with humidity; sea
salt lowers glass surface resistance by orders of magnitude and raises leakage current) from the \emph{assumed model
representation} (the logistic form, its width, $\mathit{RH}_{50}$ and $\Delta_s=15\%$~RH, none of which has been fitted
to data).
\end{response}

\begin{comment}{\S11, soiling}
State that the absolute soiling losses are scenario estimates.
\end{comment}
\begin{response}
Added to Section~IV-B: ``The absolute soiling losses are scenario estimates, because the salt and dust deposition rates
were not measured at the reference site.''
\end{response}

\begin{comment}{\S13 and \S22.4, economic uncertainty}
State that the economic conclusions are location- and year-dependent; add sensitivity to the discount rate (6--12\%) and
O\&M escalation (2--7\%).
\end{comment}
\begin{response}
Both are done:
\begin{itemize}[nosep]
\item New Table~XI gives the LCOEs of the three plants for discount rates of 6, 9 and 12\% and O\&M escalation of 2, 5
and 7\%/yr. The financing assumptions shift all LCOEs but not their ranking: the gap to land PV ranges from
\FinGapFMin{} to \FinGapFMax~INR/kWh for freshwater FPV and from \FinGapSMin{} to \FinGapSMax~INR/kWh for saline FPV.
\item Both parameters are now sampled in the Monte Carlo. Freshwater FPV is more expensive than land PV in \PLcoeF\% of
samples, and saline FPV in \PLcoeS\%.
\item Sections~III-H and IV-I state that the economic results are location- and year-dependent and specific to the
stated Indian-market assumptions.
\end{itemize}
\end{response}
\change{Sections~III-H, III-I, IV-I; Table~XI (new); Table~XIII.}

\begin{comment}{\S14 and \S22.2, baseline vs.\ Monte Carlo}
Label baseline and Monte Carlo results consistently, e.g.\ saline baseline $-4.3$\% vs.\ median $-4.94$\%.
\end{comment}
\begin{response}
Done throughout. The abstract now reports the baseline differences and then the Monte Carlo medians with their 95\%
intervals. Section~IV-H opens with: ``Baseline values and Monte Carlo medians answer different questions and are
reported separately.'' It continues: ``For freshwater FPV, the baseline 25-year difference is $+\GainLifeF$\%, while the
Monte Carlo median is $+$\MCLifeMedF\% with a 95\% interval of [\MCLifeLoF, \MCLifeHiF]\%.'' The saline and mitigated
cases follow in the same form. The section also explains why medians and baselines differ: several ranges are asymmetric
about the baseline, notably the FPV heat-loss coefficients. The conclusion uses the same two-part structure.
\end{response}

\begin{comment}{\S15--16, interpretation}
Emphasize that the freshwater advantage is conditional rather than robust; phrase the saline result as ``the model
predicts a lifetime penalty for saline FPV across most of the plausible parameter space''.
\end{comment}
\begin{response}
Both phrasings are adopted verbatim in the abstract, Section~IV-H and the conclusion.
\end{response}

\begin{comment}{\S17, $\Delta T_d^*$}
Turn the 2.7~K threshold into a practical recommendation.
\end{comment}
\begin{response}
Section~V-B now states: ``the dew-point (humidity) increment over the water should be measured before a lifetime benefit
is claimed, because the benefit vanishes above about \DTdStar~K.'' The final paragraph of the conclusion repeats it as a
recommendation for FPV assessments.
\end{response}

\begin{comment}{\S18, Fig.~12 in the Discussion}
Refer to the Monte Carlo/Spearman analysis more prominently in the Discussion.
\end{comment}
\begin{response}
New Section~V-B (``Decision-Relevant Drivers'') builds on Fig.~12(b). It identifies which inputs decide each comparison
and should therefore be measured first at a candidate site. The validation protocol (Section~V-E) refers back to it.
\end{response}

\begin{comment}{\S21, novelty claim}
Soften ``To the author's knowledge, no study yet resolves\ldots''.
\end{comment}
\begin{response}
Changed to: ``Among the studies reviewed here, none was identified that resolves, within one consistent microclimate,
\ldots''.
\end{response}

\begin{comment}{\S22.3, Monte Carlo parameter table}
Provide a table with parameter, baseline, distribution, bounds and justification/reference.
\end{comment}
\begin{response}
Added as Appendix~A, Table~XIII, covering all \NParams{} parameters. The table is generated from the same parameter
definitions that drive the Monte Carlo (\texttt{fpv\_analysis/uncertainty.py}), so it cannot drift from the code. The
justifications cite measured FPV heat-loss coefficients, ISO~9223 chloride classes and open-rack defaults where
available. All other entries are labelled ``Assumption'' or ``Calibration constant'' where no source exists.
\end{response}
\change{Appendix~A and Table~XIII (new); Section~III-I.}

\begin{comment}{\S22.1, model vs.\ validation wording}
Use ``model prediction'' rather than ``performance'' where appropriate.
\end{comment}
\begin{response}
The abstract and conclusion now say ``the model predicts'' throughout. Results remain prefaced by the statement that all
results are model predictions under the baseline assumptions.
\end{response}

\begin{comment}{\S23, minor issues}
Year-1 vs.\ year 1; notation; symbol definitions; salt vs.\ chloride deposition; units; capitalization; references.
\end{comment}
\begin{response}
\begin{itemize}[nosep]
\item ``Year-1'' is used as an adjective (``year-1 yield'') and ``year 1'' as a noun, consistently.
\item $R$, $R_{\mathrm{TH}}$, $R_{\mathrm{PID}}$ and $R_{\mathrm{C}}$ are typeset uniformly. The superscript
``ref'' denotes the reference rates.
\item Section~III-D now defines the two salt quantities explicitly. The \emph{salt deposition} $r_s$ is the sea-salt
mass retained on the glass (g/m$^2$/day), which drives soiling and leakage. The \emph{chloride deposition}
$D_{\mathrm{Cl}}$ is the chloride-ion deposition rate of the atmosphere in the ISO~9223 sense (mg/m$^2$/day), which
drives corrosion.
\item Units of deposition are given consistently in the text and tables.
\item ``FPV'' is capitalized throughout.
\item References use the IEEEtran style. Volume, page and DOI details for the references added in the first revision are
being completed from the publisher records.
\end{itemize}
\end{response}

\section*{Remaining limitation}
The absence of measured FPV and site-weather data remains the main limitation. The manuscript now presents the work as
a physically informed scenario and uncertainty framework that requires field calibration, and it proposes the
measurement protocol needed to do so.

\end{document}
